\section{Magnetized plasmas}
\label{sec:mag}

A magnetic field confines particles to gyro-orbits of radius $\rho$ across
the field while leaving them free along it. This introduces a second small
parameter, $\rho/L$, that works even without collisions. The closure
problem becomes anisotropic: the perpendicular direction is closed by
gyromotion, and the parallel direction keeps all the difficulties of
Sections~\ref{sec:plasma} and~\ref{sec:collisionless}.

\subsection{Collisionless, across the field: double-adiabatic laws}

When the fields vary slowly compared with the gyration, each particle
conserves its magnetic moment $m v_\perp^2/2B$, and the pressure tensor is
gyrotropic, $\mathbf P=p_\perp(\mathbf I-\bb\bb)+p_\parallel\bb\bb$ with
$\bb=\bB/B$. If the heat fluxes along the field are neglected, the two
pressures obey the double-adiabatic laws of Chew, Goldberger and Low
\cite{CGL56},
\begin{equation}
  \frac{\dd}{\dd t}\!\left(\frac{p_\perp}{nB}\right)=0,
  \qquad
  \frac{\dd}{\dd t}\!\left(\frac{p_\parallel B^2}{n^3}\right)=0 .
\end{equation}
The first is conservation of the magnetic moment; the second is
conservation of the parallel action for a flux tube whose length scales as
$B/n$. This is an algebraic closure with a real small parameter in the
perpendicular direction and none in the parallel direction: setting the
parallel heat fluxes to zero is the truncation of Section~\ref{sec:moment},
and it omits Landau damping. Landau-fluid models restore it by closing the
parallel heat fluxes as in Section~\ref{sec:collisionless}
\cite{Snyder97}.

\subsection{Collisional: Braginskii}

With collisions and a magnetic field the heat flux has three parts,
\begin{equation}
  \bq=-\kappa_\parallel\nabla_\parallel T-\kappa_\perp\nabla_\perp T
      \mp\kappa_\wedge\,\bb\times\nabla T ,
\end{equation}
and for electrons with $Z=1$ and $\Omega_e\tau_e\gg1$ \cite{Braginskii65}
\[
  \kappa_\parallel=3.16\,\frac{nT_e\tau_e}{m_e},\qquad
  \kappa_\perp=4.66\,\frac{nT_e}{m_e\Omega_e^2\tau_e},\qquad
  \kappa_\wedge=\frac52\,\frac{nT_e}{m_e\Omega_e},
\]
with $\Omega_e$ the gyrofrequency. The parallel coefficient is the
unmagnetized one. The perpendicular coefficient is smaller by
$(\Omega\tau)^{-2}$, many orders of magnitude in a fusion plasma, because
the step of the random walk is the gyroradius and not the mean free path.
The third term is perpendicular to the gradient and does not depend on the
collision time. It is not dissipative. The stress tensor has the same
structure, with a non-dissipative \emph{gyroviscous} part.

For a closure, the consequences are two. The parallel direction carries
almost all the heat, so errors in the parallel closure and errors in the
numerical alignment with the field dominate. And the non-dissipative terms
are closures with a sign structure (they do no work) that a learned model
must preserve.

\subsection{Closing in phase space: drift kinetics and gyrokinetics}

The same small parameter can be used before taking moments. Averaging the
kinetic equation over the gyrophase removes one velocity coordinate and the
fast gyration time scale. In gyrokinetics the ordering is
\[
  \frac{\omega}{\Omega}\sim\frac{\rho}{L}\sim\frac{\delta f}{f}
  \sim\frac{k_\parallel}{k_\perp}\ll1,\qquad k_\perp\rho\sim1,
\]
and the result is a kinetic equation in five phase-space dimensions for
rings of charge \cite{BrizardHahm07}. This is a closure in the sense of
Section~\ref{sec:problem}: the retained state is the gyrophase-averaged
distribution, the discarded information is the gyrophase dependence, and
the limiting process is the ordering above. It is asymptotic, with a
conserved adiabatic invariant behind it. Gyrofluid models are moment
closures of gyrokinetics \cite{DorlandHammett93}.

\subsection{Fluid models as a ladder of closures}

Single-fluid models of a magnetized plasma differ in what they close and
how.
\begin{center}
\renewcommand{\arraystretch}{1.3}
\begin{tabularx}{\textwidth}{@{}P{3.3cm}P{5.0cm}Y@{}}
\toprule
Model & Closure of the pressure & Closure of the electric field \\
\midrule
Ideal MHD & scalar, adiabatic: $\dd(p/n^{5/3})/\dd t=0$
  & $\bE+\bu\times\bB=0$ \\
Resistive MHD & scalar, with heating
  & $\bE+\bu\times\bB=\eta\,\mathbf J$ \\
Hall and extended MHD & scalar or Braginskii
  & generalized Ohm's law \\
Double-adiabatic & $p_\perp$, $p_\parallel$, no heat flux & ideal \\
Landau fluid & $p_\perp$, $p_\parallel$, nonlocal parallel heat flux
  & ideal or generalized \\
\bottomrule
\end{tabularx}
\end{center}
The generalized Ohm's law is the electron momentum equation solved for the
electric field,
\begin{equation}
  \bE+\bu\times\bB=\eta\,\mathbf J+\frac{\mathbf J\times\bB}{ne}
   -\frac{\nabla\cdot\mathbf P_e}{ne}
   +\frac{m_e}{ne^2}\frac{\dd\mathbf J}{\dd t},
\end{equation}
and each MHD variant keeps a subset of the terms on the right. The
electron pressure tensor in it is unclosed. In collisionless reconnection
its off-diagonal part is what breaks the frozen-in condition, which is why
ten-moment models with a heat-flux closure are used there \cite{Wang15}.
